% ============================================================================
% Chapter 12: Light Nuclei Patterns (A <= 10)
% ============================================================================
% Source: M6_Li6_Be8_ANALYSIS.md, M6_HELIUM4_ANALYSIS.md
% Status: FILLED from SoT
% Contamination check: PASSED (no banned terms)
% Contamination guard: 3D-model terminology routed to Appendix X/Q.
% Layer A text uses EDC-native vocabulary only.
% ============================================================================

\chapter{Light Nuclei Patterns}
\label{ch:light_clusters}

\source{M6\_Li6\_Be8\_ANALYSIS.md}

% ----------------------------------------------------------------------------
\section*{Abstract}
% ----------------------------------------------------------------------------

This chapter extends the pinning-cluster framework of Chapters~\ref{ch:deuterium}
and~\ref{ch:helium4} to junction networks with $A = 2 \ldots 10$ constituents.
We establish a taxonomy based on \emph{closure status}: open chains (A=2,3),
the unique closed-4 unit (A=4), and composite structures (A=5--10) built from
closed-4 cores with appendage junctions. The key finding is that the closed-4
closed-4 cluster is \emph{topologically optimal}---larger clusters either decompose
into multiple closed-4 units or suffer stability penalties from incomplete
closure. The A=8 instability (two separate closed-4 units preferred over a
single cubic cluster) emerges as a concrete prediction. This chapter does
\emph{not} provide quantitative fits for A $>$ 6; rather, it identifies
patterns and flags what would be needed to rigorize them.

% ============================================================================
\section{Minimal Cluster Taxonomy}
\label{sec:light:taxonomy}
% ============================================================================

\subsection{What ``$A$'' Means in Book IV}

In EDC terminology, $A$ is the \textbf{cluster size} of a pinned junction network:
the number of Y-junction constituents (anchor + metastable) bound together by
pinning bonds in an M$_6$ coordination lattice.

\begin{definition}[Cluster Size $A$]
    The \textbf{cluster size} $A$ counts the total number of junction states
    (both anchor and metastable) participating in a single connected pinning
    network.
\end{definition}

\subsection{Closure Classification}

Junction networks fall into three categories based on topological closure:

\begin{center}
\begin{tabular}{lp{8cm}c}
    \toprule
    Category & Definition & Example \\
    \midrule
    \textbf{Open chain} &
    Linear or branched network with distinct endpoints; no closed loop
    in the junction graph; boundary costs present &
    A=2 (deuterium) \\
    \textbf{Partial closure} &
    Network contains loops but not all internal fluxes cancel;
    some boundary cost remains &
    A=3 (triangular) \\
    \textbf{Complete closure} &
    Every junction connected to all others (complete graph);
    all internal fluxes cancel; no boundary &
    A=4 (closed-4) \\
    \bottomrule
\end{tabular}
\end{center}

\tagP

\subsection{Configuration Glossary}

Table~\ref{tab:config_glossary} provides EDC-native terminology for common
configurations.

\begin{table}[h]
\centering
\caption{Configuration types and qualitative stability expectations.}
\label{tab:config_glossary}
\begin{tabular}{llcc}
    \toprule
    $A$ & Configuration Type & Closure Status & Stability \\
    \midrule
    2 & Open pair (line segment) & None & Medium [P] \\
    3 & Triangular (partial loop) & Partial & Medium [P] \\
    4 & Closed-4 cluster & \textbf{Complete} & \textbf{High} [Dc] \\
    5 & Closed-4 + appendage & Mixed & Medium [P] \\
    6 & Closed-4 + open-2 & Mixed & Medium [P] \\
    7 & Closed-4 + triangular & Mixed & Medium [P] \\
    8 & Dual closed-4 (separated) & Complete (local) & \textbf{Unstable as single} [P] \\
    9 & Dual closed-4 + appendage & Mixed & Low--Medium [P] \\
    10 & Mixed configurations & Various & Medium [P] \\
    \bottomrule
\end{tabular}
\end{table}

% ============================================================================
\section{The A=3 Cases}
\label{sec:light:a3}
% ============================================================================

\subsection{Two Candidate Topologies}

For three junctions, two natural arrangements exist:

\paragraph{(i) Open/Partial Loop (Triangular Chain)}

Three junctions at vertices of a triangle:
\begin{itemize}
    \item 3 vertices, 3 edges (cycle graph $C_3$)
    \item Each junction has 2 neighbors (not complete)
    \item Forms a closed loop in graph sense, but \emph{not} a complete graph
    \item Some flux cancellation possible, but incomplete
\end{itemize}

\paragraph{(ii) Near-Closure but Frustrated Coordination}

Three junctions attempting closed-4 completion:
\begin{itemize}
    \item Would require 4th junction to complete closed-4 unit
    \item With only 3, one ``face'' is open
    \item Boundary cost from missing vertex
\end{itemize}

\subsection{Why A=3 Cannot Reach Closed-Unit Effect}

The A=4 closed unit achieves complete flux cancellation because:
\begin{enumerate}
    \item All 4 junctions connected to all others (6 edges)
    \item Complete graph $K_4$ has no ``outside''
    \item Net topological charge = 0
\end{enumerate}

For A=3, the graph $K_3$ (triangle) is complete, but:
\begin{itemize}
    \item Only 3 edges vs.\ 6 for A=4
    \item Cannot achieve the same localization sharing (volume too small)
    \item Flux cancellation is partial---not all modes cancel
\end{itemize}

\begin{resultbox}
    \textbf{A=3 limitation:} Triangular networks achieve partial closure
    but lack the volume and edge count for complete localization sharing.
    Binding is intermediate: $B_3 \approx 8\MeV$ vs.\ $B_4 \approx 28\MeV$.
    \tagP
\end{resultbox}

\subsection{Baseline Values for A=3}

\begin{center}
\begin{tabular}{lcc}
    \toprule
    System & $B$ (MeV) & Tag \\
    \midrule
    He-3 (2 anchor + 1 metastable) & 7.72 & [BL] \\
    H-3 (1 anchor + 2 metastable) & 8.48 & [BL] \\
    \bottomrule
\end{tabular}
\end{center}

The asymmetry (H-3 more bound than He-3) reflects the charge imbalance
affecting localization details. \tagP

\begin{edcopenbox}[title=Open: A=3 Derivation]
    \begin{itemize}
        \item Derive $B_3$ from 5D action with triangular boundary conditions
        \item Explain the He-3 / H-3 binding asymmetry from EDC first principles
        \item Quantify partial flux cancellation for 3-junction loops
    \end{itemize}
    \tagOpen
\end{edcopenbox}

% ============================================================================
\section{The A=4 Uniqueness}
\label{sec:light:a4_recap}
% ============================================================================

Chapter~\ref{ch:helium4} established that A=4 is the \emph{minimal closed
topological unit}. We summarize the key points without repetition:

\begin{enumerate}
    \item \textbf{Complete graph $K_4$:} 4 junctions, 6 edges, every pair
          connected
    \item \textbf{Closed-4 geometry:} most compact arrangement for 4 points
    \item \textbf{Complete closure:} internal fluxes cancel, net $\Phi = 0$
    \item \textbf{Maximum localization sharing:} all 4 share a common volume
    \item \textbf{Binding:} $B_4 \approx 29\MeV$ from 4-term budget
          (localization $\sim 21\MeV$, pinning $\sim 5\MeV$, surface $\sim 2\MeV$,
          closure $\sim 2\MeV$)
\end{enumerate}

\begin{resultbox}
    \textbf{The jump:} $B_4/A \approx 7.07\MeV$ vs.\ $B_3/A \approx 2.7\MeV$.
    This factor-of-2.6 jump reflects the transition from partial to complete
    topological closure.
    \tagDc + \tagP
\end{resultbox}

No larger $A$ achieves a \emph{single} closed unit with the efficiency of
A=4. Larger clusters either:
\begin{itemize}
    \item Decompose into multiple closed-4 units (preferred for A=8, 12, 16, ...)
    \item Accept incomplete closure (A=5, 6, 7, 9, 10)
\end{itemize}

\tagP

% ============================================================================
\section{A=5 to A=10: Pattern-Building by Composition}
\label{sec:light:composition}
% ============================================================================

\subsection{The Composition Principle}

For $A > 4$, the natural building strategy is:

\begin{derivationbox}
    \textbf{Closed-4 Core + Appendages:}
    \begin{equation}
        \text{Cluster}(A) = \text{Closed-4 unit} + (A - 4) \text{ appendage junctions}
        \label{eq:composition}
    \end{equation}
    The binding is approximately:
    \begin{equation}
        B(A) \approx B_4 + \sum_{\text{appendages}} N_{\text{bonds}} \times K
        + \Delta E_{\text{loc,extra}}
        \label{eq:composition_be}
    \end{equation}
    \tagP + \tagI
\end{derivationbox}

The closed-4 core provides the dominant binding ($\sim 28\MeV$), and appendage
junctions contribute incrementally through pinning bonds to the core.

\subsection{A=5: Closed-4 + 1 Appendage}

\begin{center}
\begin{verbatim}
     Closed-4 core          +    1 appendage

        o === o                    |
        |  x  |         +          o
        o === o

     4 junctions, 6 bonds      1 junction, ~3 bonds to core
\end{verbatim}
\end{center}

The 5th junction attaches to 3 vertices of the closed-4 core (forming a
capped closed-4 configuration). This adds:
\begin{itemize}
    \item 3 new edges (bonds): $\Delta E_{\text{pin}} \sim 3K \approx 2.2\MeV$
    \item Partial localization sharing: $\Delta E_{\text{loc}} \sim 2\MeV$ [P]
    \item Surface reduction: small [P]
\end{itemize}

\textbf{Estimate:} $B_5 \approx 28 + 4 = 32\MeV$. \tagP + \tagI

\subsection{A=6: Closed-4 + Open-2 (Deuteron Appendage)}

The A=6 system can be viewed as:
\begin{equation}
    \text{A=6} = \text{Closed-4 core} + \text{Open-2 pair}
\end{equation}

where the open-2 pair (like deuterium) attaches to the closed-4 core.

\begin{center}
\begin{verbatim}
     Closed-4              +      Open-2

        o === o                  o --- o
        |  x  |         +
        o === o                 (2 junctions, ~3 internal bonds)

                         Interface: ~2 bonds connecting pair to core
\end{verbatim}
\end{center}

From the source analysis:
\begin{equation}
    B_6 \approx B_4 + B_2 + B_{\text{interface}}
    \approx 28.3 + 2.2 + 1.6 \approx 32.1\MeV
    \label{eq:li6_cluster}
\end{equation}

\textbf{Baseline:} $B_6^{\text{exp}} = 32.0\MeV$ [BL]

\begin{resultbox}
    \textbf{A=6 cluster model:} Closed-4 + Open-2 with 2-bond interface gives
    $B_6 \approx 32\MeV$, matching baseline within 0.3\%.
    \tagDc + \tagP
\end{resultbox}

\subsection{A=7: Closed-4 + Triangular Appendage}

\begin{equation}
    \text{A=7} = \text{Closed-4 core} + \text{A=3 triangular appendage}
\end{equation}

The triangular appendage (3 junctions) connects to the closed-4 core,
sharing some edges. Binding estimate:
\begin{equation}
    B_7 \approx B_4 + B_3 + B_{\text{interface}} \approx 28 + 8 + 3 \approx 39\MeV
\end{equation}

\textbf{Baseline:} $B_7^{\text{exp}} = 39.24\MeV$ [BL] (for the stable A=7 isotope)

Match within 1\%. \tagP + \tagI

\subsection{Schematic: A=5, 6, 7 Composition}

\begin{center}
\begin{tabular}{cccc}
    \toprule
    $A$ & Structure & Model $B$ & Baseline $B$ \\
    \midrule
    5 & Closed-4 + 1 & $\sim 32\MeV$ [I] & (pending) \\
    6 & Closed-4 + Open-2 & $32.1\MeV$ [Dc] & $32.0\MeV$ [BL] \\
    7 & Closed-4 + Triangular & $\sim 39\MeV$ [I] & $39.24\MeV$ [BL] \\
    \bottomrule
\end{tabular}
\end{center}

\begin{edcopenbox}[title=Open: Composition Rigor]
    \begin{itemize}
        \item Derive interface bond counting from 5D geometry
        \item Quantify localization sharing between core and appendages
        \item Explain why closed-4 + appendages is preferred over
              ``larger closed unit''
    \end{itemize}
    \tagOpen
\end{edcopenbox}

% ============================================================================
\section{The A=8 Instability Case}
\label{sec:light:a8}
% ============================================================================

\subsection{The Puzzle}

For A=8, two geometries compete:

\begin{enumerate}
    \item \textbf{Single cubic cluster:} 8 junctions at vertices of a cube
          (8 vertices, 12 edges, each vertex has 3 neighbors)
    \item \textbf{Dual closed-4 units:} Two separate closed-4 clusters
          (2 $\times$ 4 junctions, each closed)
\end{enumerate}

\textbf{Observation:} The A=8 system with 4 anchor + 4 metastable junctions
is \emph{unstable}---it spontaneously separates into two closed-4 units.

\textbf{Baseline:}
\begin{center}
\begin{tabular}{lcc}
    \toprule
    Configuration & $B$ (MeV) & Status \\
    \midrule
    A=8 (single cluster) & 56.50 & Unstable [BL] \\
    $2 \times$ Closed-4 & $2 \times 28.3 = 56.6$ & Stable [BL] \\
    \midrule
    \textbf{Difference} & $-0.09$ & Favors separation \\
    \bottomrule
\end{tabular}
\end{center}

\subsection{EDC Explanation 1: Closure Saturation}

\begin{derivationbox}
    \textbf{Hypothesis A:} Each closed-4 unit achieves \emph{complete}
    flux cancellation internally. When two such units approach, their
    closures are already saturated---merging into a single 8-junction
    cluster would \emph{reopen} the topology.

    The cube has 6 faces (open boundaries) vs.\ two closed-4 units with
    no faces (closed). Opening the topology costs energy.
    \tagP
\end{derivationbox}

\subsection{EDC Explanation 2: Localization Geometry}

\begin{derivationbox}
    \textbf{Hypothesis B:} The localization energy depends on volume
    per junction. In the closed-4 unit:
    \begin{equation}
        V_{\text{tet}}/4 \approx 0.08 L_0^3 \text{ per junction}
    \end{equation}
    In the cube:
    \begin{equation}
        V_{\text{cube}}/8 \approx 1.0 L_0^3 \text{ per junction}
    \end{equation}
    The cube is \textbf{12$\times$ less volume-efficient}, meaning
    less localization energy per junction.
    \tagP
\end{derivationbox}

\subsection{Quantitative Estimate}

From the source analysis:
\begin{itemize}
    \item Cube (A=8): $B \approx 55\MeV$ [I]
    \item Dual closed-4: $B = 56.6\MeV$ [Dc]
    \item Difference: $\sim 1.6\MeV$ in favor of separation
\end{itemize}

The observed difference is $0.09\MeV$---same sign, but the model
overshoots by a factor of $\sim 20$. This indicates:
\begin{enumerate}
    \item The instability prediction is \textbf{qualitatively correct}
    \item The magnitude requires refined geometry calculations
\end{enumerate}

\begin{resultbox}
    \textbf{A=8 instability:} The pinning model correctly predicts that
    a single 8-junction cluster is less stable than two separated
    closed-4 units. The sign is correct; the magnitude is approximate.
    \tagP + \tagI
\end{resultbox}

\begin{edcopenbox}[title=Open: A=8 Precision]
    What calculation would decide between Hypothesis A (closure saturation)
    and Hypothesis B (localization geometry)?
    \begin{itemize}
        \item Compute flux energy for cube vs.\ 2$\times$closed-4 units from 5D action
        \item Determine interface energy if two closed-4 units touch but don't merge
        \item Refine localization formula for 8-junction configurations
    \end{itemize}
    \tagOpen
\end{edcopenbox}

% ============================================================================
\section{Tables}
\label{sec:light:tables}
% ============================================================================

\subsection{Table 12.1: Baseline Binding Energies}

\begin{table}[h]
\centering
\caption{Baseline binding energies for $A = 2 \ldots 10$ systems.}
\label{tab:baselines}
\begin{tabular}{cccccc}
    \toprule
    $A$ & System & $B$ (MeV) & $B/A$ (MeV) & Stability & Tag \\
    \midrule
    2 & Deuterium & 2.224 & 1.11 & Stable & [BL] \\
    3 & He-3 & 7.72 & 2.57 & Stable & [BL] \\
    3 & H-3 & 8.48 & 2.83 & Stable & [BL] \\
    4 & He-4 & 28.30 & 7.07 & Stable & [BL] \\
    5 & (various) & --- & --- & --- & Pending import \\
    6 & Li-6 & 32.0 & 5.33 & Stable & [BL] \\
    7 & Li-7 & 39.24 & 5.61 & Stable & [BL] \\
    8 & Be-8 & 56.50 & 7.06 & \textbf{Unstable} & [BL] \\
    8 & 2$\times$He-4 & 56.59 & 7.07 & Reference & [BL] \\
    9 & (various) & --- & --- & --- & Pending import \\
    10 & (various) & --- & --- & --- & Pending import \\
    \bottomrule
\end{tabular}
\end{table}

\begin{edcopenbox}[title=Open: Baseline Import]
    Import complete baseline dataset for $A = 5, 9, 10$ into Appendix~Q
    (non-binding data table).
    \tagOpen
\end{edcopenbox}

\subsection{Table 12.2: EDC Pattern Classification}

\begin{table}[h]
\centering
\caption{EDC pattern classification for $A = 2 \ldots 10$.}
\label{tab:pattern}
\begin{tabular}{cllcc}
    \toprule
    $A$ & Topology Class & Closure & Stability & Tag \\
    \midrule
    2 & Open pair & None & Medium & [Dc] \\
    3 & Triangular loop & Partial & Medium & [P] \\
    4 & Complete closed-4 & \textbf{Full} & \textbf{High} & [Dc] \\
    5 & Closed-4 + 1 appendage & Mixed & Medium & [P] \\
    6 & Closed-4 + open-2 & Mixed & Medium & [Dc] \\
    7 & Closed-4 + triangular & Mixed & Medium & [P] \\
    8 (single) & Cubic cluster & None & \textbf{Low} & [P] \\
    8 (dual) & 2$\times$ closed-4 & Full (local) & High & [Dc] \\
    9 & 2$\times$closed-4 + 1 & Mixed & Medium & [P] \\
    10 & 2$\times$closed-4 + open-2 & Mixed & Medium & [P] \\
    \bottomrule
\end{tabular}
\end{table}

\subsection{Table 12.3: Sensitivity Analysis}

\begin{table}[h]
\centering
\caption{Qualitative sensitivity of stability ranking to parameter variations.}
\label{tab:light:sensitivity}
\begin{tabular}{lccc}
    \toprule
    Variation & A=4 rank & A=6 rank & A=8 stability \\
    \midrule
    Central values & Highest & High & Unstable (single) \\
    $K + 10\%$ & Highest & High & Unstable \\
    $K - 10\%$ & Highest & High & Unstable \\
    Closure term ON & Highest & High & Unstable \\
    Closure term OFF & High (reduced) & Medium & Marginal \\
    \bottomrule
\end{tabular}
\end{table}

\textbf{Interpretation:} The A=4 ``highest stability'' ranking is robust
to $\pm 10\%$ variations in $K$. The A=8 instability (single cluster)
persists across all variations, though it becomes marginal if the
closure term is turned off. \tagP

\subsection{Table 12.4: Epistemic Ledger}

\begin{table}[h]
\centering
\caption{Epistemic status of Chapter~\ref{ch:light_clusters} claims.}
\label{tab:light:epistemic}
\begin{tabular}{p{6cm}cc}
    \toprule
    Claim & Tag & Confidence \\
    \midrule
    Baseline values ($B_2, B_3, B_4, B_6, B_7, B_8$) & [BL] & Established \\
    A=4 is unique closed unit & [Dc] & High \\
    A=3 has partial closure only & [P] & Medium \\
    A=6 = closed-4 + open-2 structure & [Dc] & High \\
    A=8 single cluster unstable & [P] & High (sign) \\
    Composition principle (Eq.~\ref{eq:composition}) & [P]+[I] & Medium \\
    Volume efficiency argument & [P] & Medium \\
    Closure saturation hypothesis & [P] & Low--Medium \\
    Quantitative A=8 difference ($1.6\MeV$) & [I] & Low (factor $\sim 20$ off) \\
    \bottomrule
\end{tabular}
\end{table}

% ============================================================================
\section{Falsifiability Hooks}
\label{sec:light:falsify}
% ============================================================================

\begin{edcopenbox}[title=Falsifiability Hooks]
    \begin{enumerate}
        \item \textbf{A=4 non-uniqueness:} If another $A \leq 10$ achieves
              $B/A > 7.07\MeV$, the closed-4 uniqueness claim fails.

        \item \textbf{A=8 stability:} If the A=8 single-cluster configuration
              is ever observed as metastable with lifetime $> 10^{-10}$ s,
              the instability prediction is wrong.

        \item \textbf{A=6 structure:} If the A=6 system is shown to have
              a geometry inconsistent with ``closed-4 + open-2,'' the
              composition principle fails.

        \item \textbf{Closure-term sensitivity:} If removing the closure
              contribution from $B_4$ still predicts highest stability for
              A=4, then closure is not the dominant mechanism.

        \item \textbf{A=3 vs.\ A=4 gap:} If the binding gap between A=3
              and A=4 can be explained without invoking topological closure,
              an alternative mechanism exists.

        \item \textbf{A=10 pattern:} If A=10 is more stable per junction than
              $2 \times$ A=5 or A=4 + A=6, the composition principle breaks.
    \end{enumerate}
\end{edcopenbox}

% ============================================================================
\section{Minimal Closed-4 Unit and Emission Bias}
\label{sec:light:emission}
% ============================================================================

\subsection{Minimal Closed Unit (A=4)}

Within Book IV, the smallest pinned-junction cluster that supports complete
closure/cancellation while remaining compatible with uniform coordination
constraints is the \textbf{closed-4 unit}. This unit is treated as the first
configuration where the closure term can fully ``turn on'' rather than remain
partial or frustrated. \tagP

\subsection{Operational Consequence}

Whenever a larger pinned cluster must reconfigure by releasing a self-contained
subcluster, the closed-4 unit is the minimal ``exportable'' object that:
\begin{enumerate}
    \item Preserves internal closure
    \item Minimizes boundary/interface cost per released junction
    \item Leaves a structurally simpler remainder
\end{enumerate}

This motivates a generic \textbf{emission bias toward closed-4 release}
in decay-like rearrangements of heavy coordination clusters. \tagP

\subsection{Forward Link (to Part E)}

In Part E (decay systematics), we will treat observed dominant release
channels as evidence that the system preferentially emits the minimal
closed unit, and we will test whether coordination-frustration metrics
predict where such releases become overwhelmingly favored.
\tagP

\begin{edcopenbox}[title=Open: Variational Proof]
    A first-principles proof requires an explicit Book IV variational statement:
    define the allowed configuration class and the effective energy functional
    (pinning + surface + localization + closure), then show that the closed-4
    unit is an extremum and (locally or globally) optimal among nearby competitors.
    \tagOpen
\end{edcopenbox}

% ============================================================================
% Observer-side projection
% ============================================================================

\begin{observerbox}
Observer-side projection note: quoted terms are measurement labels for the
3D/4D projection of the 5D objects defined in this chapter; they do not
imply any conventional mechanism.

\begin{itemize}
    \item Junction network (A $\leq$ 10) $\leftrightarrow$ light cluster (projection label)
    \item \ClusterState{} binding $\leftrightarrow$ measured binding energy (projection label)
\end{itemize}

What the observer records: the junction network compositions for $A \leq 10$
map to measured binding energies and stability patterns. The model scope
terminates at $A \sim 10$ where additional effects become dominant.
\end{observerbox}

% ============================================================================
\section{Summary}
\label{sec:light:summary}
% ============================================================================

\begin{resultbox}
    \textbf{What Chapter~\ref{ch:light_clusters} adds:}
    \begin{itemize}
        \item Taxonomy of junction networks by closure status
              (none / partial / complete)
        \item Explanation of A=3 as partial closure (triangular loop)
        \item Confirmation that A=4 is unique closed unit (closed-4)
        \item Composition principle: larger $A$ built from closed-4 cores
              + appendages
        \item A=8 instability prediction: single cube $<$ dual closed-4 units
        \item Qualitative match for A=6 ($32.1$ vs.\ $32.0\MeV$) and
              A=7 ($\sim 39$ vs.\ $39.24\MeV$)
    \end{itemize}

    \textbf{Key limitation:} The model is qualitatively successful but
    quantitatively approximate for $A > 4$. The A=8 instability sign
    is correct but magnitude is off by factor $\sim 20$.
\end{resultbox}

\bigskip
\noindent
\textbf{Forward references:}
\begin{itemize}
    \item Chapter~\ref{ch:barrier_release}: Tunneling systematics (release patterns)
    \item Chapter~\ref{ch:frustrated}: Coordination frustration in larger
          clusters
    \item Chapter~\ref{ch:high_coord}: High-coordination junction networks
    \item Chapter~\ref{ch:repro}: Reproducibility and data provenance
\end{itemize}

\bigskip
\noindent
\textbf{Derivation chain:}
\begin{equation*}
    B_4 \text{ (Ch.~\ref{ch:helium4})}
    \xrightarrow{\text{Composition}}
    B_6 \approx B_4 + B_2 + 2K
    \xrightarrow{\text{Pattern}}
    B_A \approx n_4 B_4 + n_{\text{app}} K
\end{equation*}

where $n_4$ is the number of closed-4 cores and $n_{\text{app}}$ is the
number of appendage bonds.

